\chapter{Testing}

Testing is how you keep change safe.
Tests should reflect how the system is used, not just how it is built.

\section{Unit tests}
Unit tests validate small pieces of logic quickly.
Use them for validation rules and pure functions.

\section{Integration tests}
Integration tests validate the system against real dependencies.
Use them for database access and external APIs.

\section{Contract tests}
Contract tests ensure client and server stay compatible.
They are valuable in multi-team environments.

\section{Test data}
Use deterministic fixtures.
Avoid shared state across tests.

\section{Test pyramid}
Balance unit, integration, and end-to-end tests.
Most tests should be fast unit tests.
Use end-to-end tests for critical user journeys only.

\section{End-to-end tests}
End-to-end tests verify the full system.
They are slow and brittle if overused.
Limit them to core flows such as sign-in and checkout.

\section{Property-based testing}
Property-based tests explore many inputs automatically.
They are effective for validators and parsers.
Use them to uncover edge cases not covered by examples.

\section{Fuzz testing}
Fuzzing sends random input to find crashes.
It is useful for parsers and boundary-heavy code.
Run fuzz tests on a schedule to find regressions.

\section{Mutation testing}
Mutation testing measures test quality by introducing small code changes.
If tests still pass, coverage is weak.
Use mutation testing sparingly due to cost.

\section{Flaky tests}
Flaky tests erode trust in the test suite.
Quarantine flaky tests and fix them quickly.
Track flakiness as a metric.

\section{Test environments}
Keep test environments consistent with production.
Use the same database engine and similar configuration.
Environment drift causes false confidence.

\section{CI pipelines}
CI should run fast tests on every commit.
Longer integration tests can run on a schedule or before release.
Fail fast for obvious regressions.

\section{Test data management}
Reset data between tests to avoid cross-test pollution.
Use database transactions or containerized databases.
Test data should be reproducible.

\section{Contract testing}
Contract tests prevent breaking changes between teams.
Store contracts in version control.
Fail builds when contracts are violated.

\section{Mocking strategy}
Mock external dependencies that are slow or unreliable.
Avoid mocking your own code.
Use real dependencies for integration tests.

\section{Performance testing}
Performance tests validate latency budgets.
Run them before major releases.
Track regressions over time.

\section{Canary testing}
Canary releases are a form of testing in production.
Use feature flags and gradual rollouts.
Monitor metrics closely during canaries.

\section{Example test matrix}
\begin{tabularx}{\textwidth}{@{}lX@{}}
\toprule
Test type & Primary goal \\
\midrule
Unit & validate logic quickly \\
Integration & validate dependencies and IO \\
End-to-end & validate user journeys \\
\bottomrule
\end{tabularx}

\section{Example test config}
\begin{lstlisting}[language=YAML]
tests:
  unit: "pytest -m unit"
  integration: "pytest -m integration"
  e2e: "pytest -m e2e"
\end{lstlisting}

\section{Snapshot testing}
Snapshot tests capture output for regression detection.
Use them for stable outputs only.
Large snapshots are hard to review and maintain.

\section{Database migration tests}
Migrations should be tested against realistic data sizes.
Run migrations on a copy of production schema when possible.
Migration failures are costly in production.

\section{Test isolation}
Tests must be isolated to be reliable.
Avoid shared global state.
Use per-test databases or transaction rollbacks.

\section{Test coverage}
Coverage is a signal, not a goal.
Focus on critical code paths and edge cases.
High coverage with weak assertions is still risky.

\section{Data factories}
Data factories create realistic test data.
Factories reduce duplication and keep tests readable.
Keep factories deterministic to avoid flakiness.

\section{Golden path testing}
Ensure at least one golden path test per major feature.
Golden paths verify the most common user flows.
They are useful as smoke tests in CI.

\section{Security testing}
Security tests validate common attack vectors.
Include tests for authorization boundaries.
Security tests are critical for high-risk endpoints.

\section{Test reports}
Test reports should be readable and actionable.
Include failure context, logs, and artifacts.
Poor reports slow down debugging.

\section{Static analysis}
Static analysis catches issues before runtime.
Use linters and type checkers where possible.
Static analysis complements tests, it does not replace them.

\section{Test selection}
Run only relevant tests when possible.
Use change detection to scope test runs.
Test selection speeds up feedback without sacrificing coverage.

\section{Parallelization}
Parallel test execution reduces CI time.
Ensure tests are isolated before parallelizing.
If tests share state, parallelization increases flakiness.

\section{Release testing}
Release testing validates the system in a staging environment.
Run a small suite of critical tests before production.
Release tests reduce deployment risk.

\section{Example CI pipeline}
\begin{lstlisting}[language=YAML]
pipeline:
  stages: [lint, unit, integration, e2e]
  lint: "ruff check ."
  unit: "pytest -m unit"
  integration: "pytest -m integration"
  e2e: "pytest -m e2e"
\end{lstlisting}

\section{Test ownership}
Assign ownership for critical tests.
Owners keep tests updated as the system evolves.
Ownership reduces abandoned or stale tests.

\section{Test strategy document}
Maintain a test strategy document per service.
It should describe test types and critical flows.
The document keeps testing aligned with architecture.

\section{Test debt}
Test debt accumulates when features ship without coverage.
Track test debt explicitly and pay it down.
High test debt increases release risk.

\section{Flaky test triage}
Create a triage path for flaky tests.
Quarantine, investigate, and fix.
Do not ignore flakiness, it spreads.

\section{Non-functional testing}
Non-functional tests include performance and reliability tests.
Run them for critical services on a schedule.
Non-functional testing prevents hidden regressions.

\section{Contract evolution}
When contracts change, update tests in lockstep.
Use versioned contracts for breaking changes.
Contract evolution should be deliberate.

\section{Example test strategy outline}
\begin{lstlisting}[language=YAML]
strategy:
  goals:
    - "fast feedback"
    - "protect critical flows"
  test_types:
    - unit
    - integration
    - e2e
\end{lstlisting}

\section{Environment parity}
Test environments should mirror production where possible.
Differences in databases or caches lead to missed bugs.
Document any unavoidable differences.

\section{Test data seeding}
Seed data should be minimal and purposeful.
Large seed sets slow tests and hide intent.
Use factories for most test data.

\section{Test containers}
Containers can provide consistent dependencies.
Use containers for databases and queues in integration tests.
Containerized tests reduce environment drift.

\section{Test stability metrics}
Track pass rate and duration trends.
Large changes in duration can indicate regressions.
Stability metrics help keep the suite healthy.

\section{Example stability metric}
\begin{lstlisting}[language=YAML]
stability:
  pass_rate: 0.995
  avg_duration_minutes: 12
\end{lstlisting}

\section{Testing in production}
Testing in production includes canaries and shadow traffic.
Shadow traffic validates behavior without user impact.
Always gate production tests with feature flags.

\section{Contract test example}
\begin{lstlisting}[language=JSON]
{
  "contract": "tasks.create",
  "request": {"title": "Test"},
  "response": {"id": "task_1", "title": "Test"}
}
\end{lstlisting}

\section{Test naming}
Clear names make failures easier to diagnose.
Include the behavior and context in the test name.
Avoid generic names like \code{test\_1}.

\section{Test refactoring}
Tests should be refactored like production code.
Remove duplication and improve clarity.
Unmaintained tests become liabilities.

\section{Observability for tests}
Collect logs and traces for failing tests.
Artifacts speed up debugging.
If a test fails without context, fix the test.

\section{Test environments table}
\begin{tabularx}{\textwidth}{@{}lX@{}}
\toprule
Environment & Purpose \\
\midrule
CI & fast feedback \\
Staging & release validation \\
Production & canary and shadow traffic \\
\bottomrule
\end{tabularx}

\section{Data migration testing}
Test migrations with realistic data volumes.
Validate both forward and rollback migrations.
Migration testing prevents production outages.

\section{Performance regression tests}
Keep a small suite of performance tests.
Track p95 and p99 latency over time.
Block releases on significant regressions.

\section{Chaos testing in staging}
Chaos tests validate system behavior under failure.
Run chaos tests in staging before major releases.
Chaos testing reveals missing resilience.

\section{Test suite hygiene}
Remove obsolete tests regularly.
Obsolete tests waste time and reduce clarity.
Hygiene keeps the suite fast and relevant.

\section{Test data cleanup}
Clean up test data after runs.
Orphaned data can pollute environments.
Cleanup keeps environments stable.

\section{Test performance budgets}
Set budgets for total test runtime.
If tests exceed the budget, optimize or split them.
Runtime budgets keep CI fast.

\section{Example test budget}
\begin{lstlisting}[language=YAML]
test_budget:
  total_minutes: 20
  unit_minutes: 5
  integration_minutes: 10
\end{lstlisting}

\section{Test prioritization}
Prioritize tests that protect critical revenue paths.
Low-value tests can run less frequently.
Prioritization keeps CI fast and focused.

\section{Flake quarantine}
If a test flakes, quarantine it and assign an owner.
Quarantined tests should not block releases.
Track quarantine counts as a health metric.

\section{Boundary testing}
Boundary tests validate behavior at limits.
Test max sizes, empty inputs, and invalid types.
Boundary tests catch edge cases before production does.

\section{Concurrency testing}
Concurrent access can break otherwise correct logic.
Test for races in caches, queues, and database updates.
Use deterministic schedulers when possible.

\section{Clock control}
Time-dependent logic needs controlled clocks in tests.
Use fake clocks to avoid sleeps and flakiness.
Clock control keeps tests fast and reliable.

\section{Test doubles}
Use the right kind of double for the situation.
Stubs provide data, mocks verify interactions, fakes provide behavior.
Choose doubles that keep tests readable and stable.

\section{Test doubles table}
\begin{tabularx}{\textwidth}{@{}lX@{}}
\toprule
Double & Purpose \\
\midrule
Stub & returns fixed data \\
Mock & verifies interactions \\
Fake & simplified behavior \\
\bottomrule
\end{tabularx}

\section{Schema testing}
APIs and events should be validated against schemas.
Schema tests catch breaking changes early.
Use versioned schemas for evolvable systems.

\section{Replay testing}
Replay tests run production traffic against new versions.
They catch behavior drift without user impact.
Ensure data is anonymized before replay.

\section{Data privacy in tests}
Never use raw production data in lower environments.
Anonymize or synthesize test data.
Privacy violations in tests are still violations.

\section{Test review}
Review tests with the same rigor as production code.
Reject tests with weak assertions or poor isolation.
Test review improves suite quality over time.

\section{Risk-based testing}
Prioritize tests based on business impact.
High-risk features need deeper coverage.
Risk-based testing keeps efforts aligned with priorities.

\section{Example risk matrix}
\begin{tabularx}{\textwidth}{@{}lX@{}}
\toprule
Risk level & Testing focus \\
\midrule
High & end-to-end plus integration \\
Medium & integration plus unit \\
Low & unit only \\
\bottomrule
\end{tabularx}

\section{API testing}
API tests should validate status codes and error shapes.
Include authentication and authorization scenarios.
APIs are the contract boundary for most systems.

\section{Queue processing tests}
Async workers need tests for retries and idempotency.
Validate that duplicates do not corrupt data.
Queue tests prevent silent failures in background work.

\section{Feature flag testing}
Test both enabled and disabled flag paths.
Ensure flags are removed after rollout completes.
Flags are code paths that can rot quickly.

\section{Rollback testing}
Rollbacks should be tested like deployments.
Validate that schema changes are backward compatible.
Rollback tests reduce risk during incidents.

\section{Load test fixtures}
Use realistic workloads and data distributions.
Synthetic uniform loads hide bottlenecks.
Load tests should match real usage patterns.

\section{Test tagging}
Tag tests by scope and criticality.
Use tags to run a minimal smoke suite quickly.
Tags help balance speed and coverage.

\section{Example test tags}
\begin{lstlisting}[language=YAML]
tags:
  smoke: ["auth", "checkout"]
  slow: ["e2e", "load"]
\end{lstlisting}

\section{Schema migration rehearsal}
Run migrations in staging with production-like volumes.
Measure duration and verify data integrity.
Rehearsals prevent surprises during release windows.

\section{Eventual consistency tests}
Systems with async flows need tests for timing windows.
Assert eventual outcomes with bounded retries.
Explicitly document expected delays.

\begin{checklistbox}{Checklist: testing}
\begin{itemize}[nosep]
  \item Keep unit tests fast.
  \item Use integration tests for critical flows.
  \item Keep fixtures deterministic.
\end{itemize}
\end{checklistbox}
